\chapter{Il sistema nervoso autonomo}

% ---------------------------------------------------------------------------
\section{Generalità}

\textbf{Sistema nervoso autonomo} (SNA): principale coordinatore delle
funzioni dell'organismo; regola la temperatura corporea e coordina le
attività degli apparati cardiovascolare, respiratorio, digerente, escretore
e riproduttivo. In questo modo controlla i livelli di acqua, elettroliti e
sostanze nutritizie disciolte nei fluidi corporei.

Il sistema nervoso centrale nel suo complesso è responsabile di: percezioni,
comportamento, memoria, riconoscimento, apprendimento, immaginazione,
ragionamento astratto, pensiero creativo.

Il SNA può essere diviso in tre parti:
\begin{itemize}
  \item il sistema nervoso simpatico (od ortosimpatico);
  \item il sistema nervoso parasimpatico;
  \item il sistema nervoso enterico (o metasimpatico).
\end{itemize}

Il SNA è formato da nervi efferenti che lasciano il SNC e innervano organi
effettori periferici: cuore, vasi, ghiandole, organi viscerali, muscolatura
liscia (\textbf{non} la muscolatura scheletrica).

\rubrica{Due neuroni in serie}
\begin{itemize}
  \item \textbf{neurone pregangliare}: corpo cellulare in un nucleo del
    tronco dell'encefalo o nella sostanza grigia del midollo spinale;
  \item \textbf{neurone postgangliare}: corpo cellulare in un ganglio
    periferico.
\end{itemize}
I due neuroni entrano in contatto sinaptico nei gangli, e l'assone del
neurone postgangliare trasmette impulsi ai muscoli involontari o alle
ghiandole.

% ---------------------------------------------------------------------------
\section{Organizzazione delle tre divisioni}

\textbf{Sistema nervoso simpatico}: origina dal midollo spinale, a livello
delle zone toraciche e lombari. I prolungamenti dei neuroni si dirigono
verso i gangli localizzati vicino al midollo spinale: si parla quindi di
\textit{fibre pregangliari} (che originano nel midollo) e di \textit{fibre
postgangliari} (che partono dal ganglio e si dirigono verso un muscolo o
una ghiandola). Solo poche fibre pregangliari contattano direttamente altri
gangli diversi da questi.

\textbf{Sistema nervoso parasimpatico}: i corpi dei neuroni si trovano
nella regione sacrale del midollo spinale e nel midollo allungato del
tronco encefalico, dove i nervi cranici III, VII, IX e X formano le fibre
pregangliari parasimpatiche. Queste, insieme a quelle che originano dal
midollo spinale, si dirigono verso gangli molto vicini all'organo da
controllare; da qui le fibre postgangliari si dirigono direttamente verso
l'organo bersaglio.

\textbf{Sistema nervoso enterico}: è formato dalle fibre nervose che
innervano i visceri.

\rubrica{Sistemi motori autonomo e somatico}
\begin{itemize}
  \item \textbf{sistema motorio somatico}: un motoneurone si estende dal
    SNC al muscolo scheletrico; gli assoni sono mielinici e conducono gli
    impulsi rapidamente;
  \item \textbf{sistema nervoso autonomo}: una catena di due motoneuroni
    (neurone pregangliare e neurone postgangliare); la conduzione è più
    lenta, dovuta ad assoni poco o per niente mielinizzati.
\end{itemize}

\rubrica{Origine anatomica delle tre divisioni}
\begin{description}
  \item[Simpatico:] corpi cellulari dei neuroni pregangliari nelle corna
    laterali dei segmenti toracici e lombari del midollo spinale. Gli
    assoni lasciano il midollo tramite le radici ventrali dei nervi
    spinali e raggiungono i corpi cellulari dei neuroni postgangliari nei
    gangli vertebrali, che formano una catena posta a lato della colonna
    vertebrale.
  \item[Parasimpatico:] corpi cellulari localizzati in diversi nuclei del
    tronco dell'encefalo e nelle corna laterali della sostanza grigia dei
    segmenti sacrali del midollo spinale. Gli assoni lasciano il SNC
    tramite nervi cranici (oculomotore, facciale, glossofaringeo e vago) e
    nervi sacrali spinali, connettendosi con i neuroni postgangliari nei
    gangli parasimpatici, distribuiti vicino agli organi innervati.
  \item[Enterico:] neuroni situati nei plessi intramurali del tratto
    gastrointestinale. Riceve impulsi dal sistema simpatico e
    parasimpatico, ma può agire autonomamente per controllare le funzioni
    motorie e secretorie dell'intestino.
\end{description}

% ---------------------------------------------------------------------------
\section{Differenze anatomiche e funzionali tra simpatico e parasimpatico}

\rubrica{Sistema nervoso simpatico}
Assoni pregangliari corti, che terminano in strutture vicine; assoni
postgangliari lunghi. Gli assoni di un singolo ganglio si distribuiscono in
maniera diffusa e innervano molti organi $\rightarrow$ le reazioni
simpatiche sono diffuse.

\rubrica{Sistema nervoso parasimpatico}
Assoni pregangliari lunghi, che terminano in gangli situati all'interno o
in vicinanza degli organi che innervano $\rightarrow$ le reazioni
parasimpatiche sono localizzate.

\rubrica{Riepilogo delle differenze anatomiche}
\begin{itemize}
  \item \textbf{lunghezza delle fibre postgangliari}: ortosimpatiche
    lunghe, parasimpatiche corte;
  \item \textbf{diramazione degli assoni}:
  \begin{itemize}
    \item assone ortosimpatico: fortemente diramato, influenza diversi
      organi;
    \item assoni parasimpatici: poche diramazioni, effetto localizzato.
  \end{itemize}
\end{itemize}

% ---------------------------------------------------------------------------
\section{La reazione di attacco o fuga}

\textbf{Ortosimpatico} (``lotta, fuga o paura''): attivato durante
l'esercizio fisico, l'eccitamento e le emergenze.

\textbf{Parasimpatico} (``riposo e digestione''): coinvolto nel risparmio
energetico.

La reazione di attacco o fuga, chiamata anche ``reazione da stress acuta''
o \textit{fight or flight response}, fu descritta da Walter Cannon nel
1920: gli animali (uomo compreso) reagiscono alle minacce con una scarica
generale del sistema nervoso simpatico. La risposta è stata
successivamente riconosciuta come la prima tappa di una sindrome generale
di adattamento allo stress e come l'individuo risponde ad esso. Cannon
capì che la risposta alle minacce è memorizzata profondamente nel nostro
cervello e rappresenta una saggezza genetica: è qualcosa di ``progettato''
per proteggerci dai danni fisici.

\rubrica{Effetti della risposta di attacco o fuga}
\begin{itemize}
  \item la frequenza respiratoria aumenta;
  \item il sangue viene deviato lontano dall'apparato digerente e diretto
    nei muscoli e in altre parti del corpo che richiedono energia
    supplementare per la corsa e la lotta;
  \item le pupille si dilatano;
  \item la consapevolezza si intensifica;
  \item la vista si acuisce;
  \item il battito cardiaco accelera;
  \item la percezione del dolore diminuisce;
  \item il sistema immunitario si attiva maggiormente;
  \item ci si prepara fisicamente e psicologicamente all'attacco o alla
    fuga;
  \item si controlla l'ambiente in cui ci si trova, ``alla ricerca del
    nemico''.
\end{itemize}

Un pericolo emotivo può affinare l'acutezza mentale, contribuendo così ad
affrontare con decisione i problemi, passando all'azione: è questo il
vantaggio di trovarsi nello stato di attacco o fuga.

% ---------------------------------------------------------------------------
\section{Organizzazione dei gangli simpatici}

\rubrica{Gangli collaterali (latero-vertebrali)}
Principali effetti prodotti dalle fibre pregangliari dirette ai gangli
collaterali:
\begin{itemize}
  \item costrizione arteriolare e riduzione di flusso ai visceri;
  \item inibizione dell'attività degli organi e delle ghiandole
    dell'apparato digerente;
  \item rilascio di glucosio dalle riserve epatiche di glicogeno;
  \item rilascio di lipidi da parte del tessuto adiposo;
  \item rilasciamento della muscolatura liscia della parete vescicale;
  \item riduzione della formazione di urina a livello dei reni;
  \item controllo di alcuni aspetti della funzione sessuale (es.\
    eiaculazione nel maschio).
\end{itemize}

\rubrica{Sostanza midollare del surrene}
Fibre pregangliari che originano tra T5 e T8. Le cellule endocrine
(neuroni gangliari specializzati) della sostanza midollare del surrene
secernono neurotrasmettitori direttamente nel circolo sanguigno: il
principale effetto è il rilascio di adrenalina e noradrenalina nella
circolazione generale.

\rubrica{Distribuzione anatomica delle fibre postgangliari simpatiche}
\begin{itemize}
  \item solo i gangli toracici e lombari superiori ricevono fibre
    pregangliari attraverso rami comunicanti bianchi;
  \item le catene gangliari cervicale, lombare inferiore e sacrale
    ricevono l'innervazione pregangliare dai segmenti toracico e lombare
    superiore mediante fibre pregangliari che ascendono o discendono
    lungo la catena;
  \item ciascun nervo spinale riceve un ramo comunicante grigio da un
    ganglio della catena del simpatico.
\end{itemize}

% ---------------------------------------------------------------------------
\section{La sezione parasimpatica}

\rubrica{Uscita craniale}
Proviene dal cervello e innerva gli organi della testa, del collo, del
torace e dell'addome. Le fibre pregangliari corrono attraverso:
\begin{itemize}
  \item il nervo oculomotore (III);
  \item il nervo facciale (VII);
  \item il nervo glossofaringeo (IX);
  \item il nervo vago (X).
\end{itemize}
I corpi cellulari dei neuroni pregangliari sono localizzati nei nuclei dei
nervi cranici del tronco encefalico.

\rubrica{Uscita sacrale}
Innerva i rimanenti organi addominali e pelvici; emerge da $S_2$--$S_4$. I
corpi cellulari pregangliari sono localizzati nella regione motoria
viscerale della sostanza grigia spinale e formano i nervi splancnici.

% ---------------------------------------------------------------------------
\section{La sezione ortosimpatica}

\rubrica{Organizzazione di base}
\begin{itemize}
  \item fuoriuscita da $T_1$--$L_2$;
  \item alimenta gli organi viscerali e le strutture delle regioni
    superficiali del corpo;
  \item contiene più gangli della divisione parasimpatica: alcuni sono
    \textit{paravertebrali}, altri \textit{prevertebrali}.
\end{itemize}

\rubrica{Gangli della catena paravertebrale}
\begin{itemize}
  \item localizzati su entrambi i lati della colonna vertebrale;
  \item collegati tramite brevi tratti della catena ortosimpatica
    paravertebrale;
  \item uniti ai rami ventrali attraverso i rami comunicanti bianchi;
  \item principali gangli: cervicale superiore, medio e inferiore, ecc.
\end{itemize}

\rubrica{Gangli prevertebrali}
\begin{itemize}
  \item spaiati, non organizzati in maniera segmentale;
  \item si riscontrano solo a livello addominale e pelvico;
  \item giacciono anteriormente alla colonna vertebrale;
  \item principali gangli: celiaco, ipogastrico superiore e inferiore.
\end{itemize}

\rubrica{Percorso delle fibre ortosimpatiche}
Indipendentemente dal bersaglio, l'origine è comune: gli assoni
pregangliari fuoriescono dal midollo spinale attraverso le radici
ventrali, entrano nei nervi spinali, li lasciano attraverso i rami
comunicanti ed entrano nella catena ortosimpatica dei gangli paravertebrali,
dove si trova il soma dei neuroni postgangliari.

Da qui esistono tre possibilità:
\begin{enumerate}
  \item formare una sinapsi con il neurone postgangliare nella catena dei
    gangli paravertebrali e ritornare ai nervi spinali per proseguire
    verso la cute;
  \item salire o scendere all'interno della catena dei gangli
    paravertebrali, formare una sinapsi con un neurone postgangliare
    all'interno della catena, e ritornare ai nervi spinali per proseguire
    verso la cute;
  \item entrare nella catena dei gangli paravertebrali, passarvi
    attraverso senza formare sinapsi, e formare il nervo splancnico che
    raggiunge gli organi toracici e addominali: questi formano una
    sinapsi a livello dei gangli prevertebrali di fronte all'aorta, da cui
    gli assoni postgangliari raggiungono gli organi.
\end{enumerate}

% ---------------------------------------------------------------------------
\section{Zone a prevalente controllo simpatico o parasimpatico}

\begin{itemize}
  \item \textbf{a prevalente controllo parasimpatico}: secrezioni
    dell'apparato gastroenterico;
  \item \textbf{a prevalente controllo ortosimpatico}: vasi, bronchi,
    utero, metabolismo;
  \item \textbf{controllati da entrambi} (con azioni opposte): cuore,
    occhio, muscolatura liscia intestinale, vescica.
\end{itemize}

Quando il sistema simpatico è attivo e il parasimpatico quiescente si ha
la risposta ``combatti o fuggi''; quando il parasimpatico è attivo e il
simpatico quiescente si ha la risposta ``riposa e digerisci''.
L'omeostasi è mantenuta grazie a uno stato di equilibrio dinamico tra le
due branche del sistema nervoso autonomo.

\rubrica{Azioni del sistema nervoso autonomo sui principali organi}
\begin{table}[htbp]
  \centering
  \small
  \renewcommand{\arraystretch}{1.3}
  \begin{tabularx}{\textwidth}{|l|X|X|}
    \hline
    \textbf{Struttura} & \textbf{Stimolazione dell'ortosimpatico} &
    \textbf{Stimolazione del parasimpatico} \\
    \hline
    Occhio (iride) & dilatazione della pupilla & costrizione della pupilla \\
    \hline
    Ghiandole salivari & riduzione della salivazione & aumento della
    salivazione \\
    \hline
    Mucosa orale & riduzione della produzione di muco & aumento della
    produzione di muco \\
    \hline
    Cuore & aumento della frequenza dei battiti e della forza di
    contrazione & diminuzione della frequenza dei battiti e della forza di
    contrazione \\
    \hline
    Polmoni & rilassamento dei bronchi & contrazione della muscolatura
    bronchiale \\
    \hline
    Stomaco & riduzione della motilità & secrezione di succo gastrico e
    aumento della motilità \\
    \hline
    Intestino tenue & riduzione della peristalsi & aumento dei processi
    digestivi \\
    \hline
    Intestino crasso & riduzione della motilità & aumento della secrezione
    e della motilità \\
    \hline
    Fegato & aumentata glicogenolisi & -- \\
    \hline
    Rene & diminuzione della diuresi & aumento della diuresi \\
    \hline
    Midollare surrenale & secrezione di adrenalina e noradrenalina & -- \\
    \hline
    Vescica & rilassamento della parete e chiusura dello sfintere &
    contrazione della parete e rilasciamento dello sfintere \\
    \hline
  \end{tabularx}
\end{table}

% ---------------------------------------------------------------------------
\section{Riepilogo sulla divisione simpatica}

\begin{itemize}
  \item la divisione simpatica del SNA comprende due catene gangliari su
    ciascun lato della colonna vertebrale, tre gangli collaterali
    anteriori al midollo spinale e la zona midollare delle due ghiandole
    surrenali;
  \item le fibre pregangliari sono relativamente brevi, poiché i gangli
    sono vicini al midollo spinale; le fibre postgangliari sono invece
    relativamente lunghe, poiché i loro organi bersaglio sono distanti da
    essi;
  \item la divisione simpatica mostra un'ampia divergenza: una singola
    fibra pregangliare può innervare fino a 32 neuroni gangliari diversi;
    ne consegue che un singolo motoneurone simpatico può controllare
    diversi effettori periferici, producendo una risposta complessa e
    coordinata;
  \item tutti i neuroni pregangliari liberano \textit{acetilcolina} a
    livello delle sinapsi con i neuroni gangliari; la maggior parte delle
    fibre postgangliari rilascia adrenalina e, in minima parte,
    acetilcolina;
  \item la risposta effettrice dipende dalla funzione del recettore di
    membrana attivato quando adrenalina e noradrenalina si legano ai
    recettori alfa o beta.
\end{itemize}

% ---------------------------------------------------------------------------
\section{Il sistema nervoso enterico (metasimpatico)}

Costituito da circa 100 milioni di neuroni (numero paragonabile a quelli
che costituiscono il midollo spinale); è per lo più indipendente dal
sistema nervoso simpatico e parasimpatico. È definito come ``cervello
addominale'' o ``cervello enterico''.

I neuroni del sistema nervoso enterico si raggruppano in due plessi:
\begin{itemize}
  \item \textbf{plesso sottomucoso}: regola l'attività secretoria del
    tubo digerente;
  \item \textbf{plesso mio-enterico}: controlla l'attività motoria
    gastrointestinale lungo tutta la sua lunghezza.
\end{itemize}

% ---------------------------------------------------------------------------
\section{Tipici riflessi regolati dal sistema nervoso autonomo}

\rubrica{Minzione (continenza)}
\begin{itemize}
  \item \textbf{muscolo detrusore}: innervato da terminazioni nervose
    sensoriali (stiramento);
  \item \textbf{muscolo sfintere interno}: muscolo liscio, innervato da
    fibre parasimpatiche (nervi pelvici) e fibre simpatiche (nervi
    ipogastrici);
  \item \textbf{muscolo sfintere esterno}: muscolo striato scheletrico,
    innervato dai motoneuroni del nervo pudendo, sotto il controllo
    volontario di centri superiori.
\end{itemize}

I segnali della branca simpatica facilitano il riempimento della vescica,
rilassando la muscolatura, inibendo l'attività dei gangli parasimpatici e
diminuendo la capacità contrattile del muscolo detrusore.
